\documentclass[10pt,a4paper]{amsart}
\usepackage{amsmath,amssymb,amsthm,graphicx}
% Retained preamble for the Stacks Project chapter morphisms.tex (Morphisms of Schemes).
% The chapter file pinned-original-source.tex carries only \input{preamble} on its first line, so no
% preamble text lives inside the chapter itself.  The Stacks Project uses one single shared preamble for
% every chapter; the snapshot stages it verbatim in sources/stacks-chow/preamble.tex.  The block below is
% copied mechanically, line for line, from that file's lines 14-18 (the Xy-pic setup) and lines 154-217
% (the native theorem declarations together with the mathematical macros).  Nothing was added, removed or
% redefined.  The wrapper class and the amsmath/amssymb/amsthm/graphicx base packages are supplied by the
% assembler, so the class/\externaldocument part of the shared preamble is not repeated here.
% For commutative diagrams we use Xy-pic
\usepackage[all]{xy}

% We use 2cell for 2-commutative diagrams.
\xyoption{2cell}
% Theorem environments.
%
\theoremstyle{plain}
\newtheorem{theorem}[subsection]{Theorem}
\newtheorem{proposition}[subsection]{Proposition}
\newtheorem{lemma}[subsection]{Lemma}

\theoremstyle{definition}
\newtheorem{definition}[subsection]{Definition}
\newtheorem{example}[subsection]{Example}
\newtheorem{exercise}[subsection]{Exercise}
\newtheorem{situation}[subsection]{Situation}

\theoremstyle{remark}
\newtheorem{remark}[subsection]{Remark}
\newtheorem{remarks}[subsection]{Remarks}

\numberwithin{equation}{subsection}

% Macros
%
\def\lim{\mathop{\mathrm{lim}}\nolimits}
\def\colim{\mathop{\mathrm{colim}}\nolimits}
\def\Spec{\mathop{\mathrm{Spec}}}
\def\Hom{\mathop{\mathrm{Hom}}\nolimits}
\def\Ext{\mathop{\mathrm{Ext}}\nolimits}
\def\SheafHom{\mathop{\mathcal{H}\!\mathit{om}}\nolimits}
\def\SheafExt{\mathop{\mathcal{E}\!\mathit{xt}}\nolimits}
\def\Sch{\mathit{Sch}}
\def\Mor{\mathop{\mathrm{Mor}}\nolimits}
\def\Ob{\mathop{\mathrm{Ob}}\nolimits}
\def\Sh{\mathop{\mathit{Sh}}\nolimits}
\def\NL{\mathop{N\!L}\nolimits}
\def\CH{\mathop{\mathrm{CH}}\nolimits}
\def\proetale{{pro\text{-}\acute{e}tale}}
\def\etale{{\acute{e}tale}}
\def\QCoh{\mathit{QCoh}}
\def\Ker{\mathop{\mathrm{Ker}}}
\def\Im{\mathop{\mathrm{Im}}}
\def\Coker{\mathop{\mathrm{Coker}}}
\def\Coim{\mathop{\mathrm{Coim}}}

% Boxtimes
%
\DeclareMathSymbol{\boxtimes}{\mathbin}{AMSa}{"02}

%
% Macros for moduli stacks/spaces
%
\def\QCohstack{\mathcal{QC}\!\mathit{oh}}
\def\Cohstack{\mathcal{C}\!\mathit{oh}}
\def\Spacesstack{\mathcal{S}\!\mathit{paces}}
\def\Quotfunctor{\mathrm{Quot}}
\def\Hilbfunctor{\mathrm{Hilb}}
\def\Curvesstack{\mathcal{C}\!\mathit{urves}}
\def\Polarizedstack{\mathcal{P}\!\mathit{olarized}}
\def\Complexesstack{\mathcal{C}\!\mathit{omplexes}}
% \Pic is the operator that assigns to X its picard group, usage \Pic(X)
% \Picardstack_{X/B} denotes the Picard stack of X over B
% \Picardfunctor_{X/B} denotes the Picard functor of X over B
\def\Pic{\mathop{\mathrm{Pic}}\nolimits}
\def\Picardstack{\mathcal{P}\!\mathit{ic}}
\def\Picardfunctor{\mathrm{Pic}}
\def\Deformationcategory{\mathcal{D}\!\mathit{ef}}

% --- appended 2026-10-12 (Chief): the project own definitions of the reference/slogan/history
% display environments, copied mechanically from the shared Stacks preamble lines 8-12
% (\usepackage{verbatim} plus the three \newenvironment lines, which the chapter preamble keeps
% verbatim from the project preamble). Nothing added, removed or redefined: these environments
% carry no mathematical content and the project itself typesets them as comments. Without these
% three lines a counted statement whose own text uses them fails with Environment slogan undefined.
\usepackage{verbatim}
\newenvironment{reference}{\comment}{\endcomment}
\newenvironment{slogan}{\comment}{\endcomment}
\newenvironment{history}{\comment}{\endcomment}
\usepackage{cleveref}
\crefname{originallemma}{original lemma}{original lemmas}
\Crefname{originallemma}{Original lemma}{Original lemmas}
\makeatletter
\newlabel{algebra-lemma-Noetherian-irreducible-components}{{00FR}{source}{}{sourcepointer-1}{}}
\newlabel{algebra-lemma-Noetherian-irreducible-components@cref}{{[originallemma][0][]00FR}{[1][0][]source}{}{}{}}
\makeatother
\crefname{originallemma}{original lemma}{original lemmas}
\Crefname{originallemma}{Original lemma}{Original lemmas}
\makeatletter
\newlabel{algebra-lemma-injective-minimal-primes-in-image}{{00FK}{source}{}{sourcepointer-2}{}}
\newlabel{algebra-lemma-injective-minimal-primes-in-image@cref}{{[originallemma][0][]00FK}{[1][0][]source}{}{}{}}
\makeatother
\crefname{originallemma}{original lemma}{original lemmas}
\Crefname{originallemma}{Original lemma}{Original lemmas}
\makeatletter
\newlabel{lemma-factor-through-point}{{0H1M}{source}{}{sourcepointer-3}{}}
\newlabel{lemma-factor-through-point@cref}{{[originallemma][0][]0H1M}{[1][0][]source}{}{}{}}
\makeatother
\title{A morphism from a reduced scheme into a set of mutually incomparable points of a locally Noetherian scheme factors through that discrete set}
\author{The Stacks Project Authors}
\date{stacks-project:morphisms.tex:local-snapshot-20261012; GNU FDL 1.2 or later}
\begin{document}
\maketitle
\noindent\textit{Editorial note.} All delivered mathematical text is the author's own, copied from the pinned Stacks Project chapter file pinned-original-source.tex; no statement or proof step was rewritten, reordered or removed.\par\medskip
\noindent\textit{Editorial note.} Following the multi-claim convention for this source, source\_url is the source-level URL https://stacks.math.columbia.edu/ and the source label lemma-factor-through-set-of-points is carried in claim\_key and here rather than in a per-claim tag URL.\par\medskip
\noindent\textit{Editorial note.} The preamble is the chapter's shared Stacks Project preamble, retained as sources/stacks-morphisms/retained-preamble.tex: a line-for-line excerpt of the shared Stacks preamble staged with this snapshot. No macro was invented or redefined, and this card adds no preamble text of its own.\par\medskip
\noindent\textit{Editorial note.} Statement span 15616-15622 and proof span 15624-15643 come from the halfA pool card operations/pool-cards/stacks-morphisms-halfA.json and were re-read in the pinned chapter file before this card was written.\par\medskip
\noindent\textit{Editorial note.} The proof's argument that the image is contained in the finite set makes explicit use of the hypothesis that there are no nontrivial specializations among the elements of E; the author's wording is preserved.\par\medskip
\noindent\textit{Editorial note.} 3 cross-references occur inside the counted spans; they are listed in external\_references as navigation pointers only.\par\medskip
\noindent\textit{Allowed standard prerequisites (editorial).}\begin{itemize}
\item The statement for a single point: a morphism from a reduced scheme into one point of a locally Noetherian scheme, with image exactly that point, factors through the inclusion of the point, cited from the same chapter. Cited at proof lines 15624-15626. Gives the singleton case on which the finite case is built.
\item The set of minimal primes of a quotient of a Noetherian ring is finite, and the minimal primes of the image of an injective ring map are contained in the image of the corresponding morphism of spectra, cited from the algebra chapter. Cited at proof lines 15636-15640. Produces a finite subset of the given set that contains all points of the image, reducing to the finite case.
\end{itemize}\medskip
\ifdefined\hypertarget\hypertarget{sourcepointer-1}{}\fi
\noindent\textit{Source reference.} Original lemma 00FR: focused-seeds/AGP/sources/stacks-snapshot-20261011/tags.txt, source lines 581--581. Same-chapter reference used in the retained author text at source line(s) 13632, 14321, 15637. It is supplied only as a navigation pointer so that the retained text compiles; the counted item delivers the author's own statement and the author's complete proof, and no part of the referenced item is copied in or substituted for the counted argument.\par\medskip
\ifdefined\hypertarget\hypertarget{sourcepointer-2}{}\fi
\noindent\textit{Source reference.} Original lemma 00FK: focused-seeds/AGP/sources/stacks-snapshot-20261011/tags.txt, source lines 574--574. Same-chapter reference used in the retained author text at source line(s) 14092, 15633. It is supplied only as a navigation pointer so that the retained text compiles; the counted item delivers the author's own statement and the author's complete proof, and no part of the referenced item is copied in or substituted for the counted argument.\par\medskip
\ifdefined\hypertarget\hypertarget{sourcepointer-3}{}\fi
\noindent\textit{Source reference.} Original lemma 0H1M: focused-seeds/AGP/sources/stacks-snapshot-20261011/tags.txt, source lines 21252--21252. Same-chapter reference used in the retained author text at source line(s) 15626. It is supplied only as a navigation pointer so that the retained text compiles; the counted item delivers the author's own statement and the author's complete proof, and no part of the referenced item is copied in or substituted for the counted argument.\par\medskip
\begin{lemma}
\label{lemma-factor-through-set-of-points}
Let $f : Y \to X$ be a morphism of schemes. Let $E \subset X$.
Assume $X$ is locally Noetherian, there are no nontrivial specializations
among the elements of $E$, $Y$ is reduced, and $f(Y) \subset E$.
Then $f$ factors through $\coprod_{x \in E} x \to X$.
\end{lemma}

\begin{proof}
When $E$ is a singleton this follows from
Lemma \ref{lemma-factor-through-point}. If $E$ is finite, then $E$
(with the induced topology of $X$) is a finite discrete space by
our assumption on specializations. Hence this case reduces to the
singleton case. In general, there is a reduction to the
case where $X$ and $Y$ are affine schemes. Say $f : Y \to X$ corresponds to
the ring map $\varphi : A \to B$. Denote $A' \subset B$ the image of $\varphi$.
Let $E' \subset \Spec(A') \subset \Spec(A)$ be the set of minimal primes
of $A'$. By Algebra, Lemma \ref{algebra-lemma-injective-minimal-primes-in-image}
the set $E'$ is contained in the image of
$\Spec(B) \to \Spec(A') \subset \Spec(A)$. We conclude that $E' \subset E$.
Since $A'$ is Noetherian we have $E'$ is finite by Algebra, Lemma
\ref{algebra-lemma-Noetherian-irreducible-components}.
Since any other point in the image of $\Spec(B) \to \Spec(A)$ is
a specialization of an element of $E'$ and in $E$, we conclude that
the image is contained in $E'$ (by our assumption on specializations between
points of $E$). Thus we reduce to the case where
$E$ is finite which we dealt with above.
\end{proof}

\end{document}
